\documentclass[lettersize,journal]{IEEEtran}
\usepackage{amsmath,amsfonts,amssymb}
\usepackage{url}
\usepackage{cite}
\usepackage{booktabs}
\usepackage{textcomp}
\usepackage{xr}
\externaldocument{main}
% \hyphenation{op-tical net-works semi-conduc-tor IEEE-Xplore}

% --- Supplementary numbering ---
\renewcommand{\theequation}{S\arabic{equation}}
\renewcommand{\thesection}{S-\Roman{section}}
\renewcommand{\thesubsection}{\thesection.\Alph{subsection}}

\begin{document}

\title{Supplementary Material for ``Graph Learning for Topology-Adaptive Fast Frequency Response in Inverter-Based Microgrids''}

\author{Tran Kim Nhat, Nguyen Van Cuong, Nguyen Phi Long, and Laurent El Ghaoui%
\thanks{Equation numbers of the form (S$x$) refer to this document; plain numbers, e.g.\ (6), refer to the main paper.}}

\maketitle

% ==================================================================
\section{Component and Frequency-Model Details}
\label{supp:models}
% ==================================================================
This section collects the device and network equations underlying the reduced-order frequency model of Sec.~II-B of the main paper; the symbols follow the main-paper Nomenclature.

\subsection{DER Device Models}
A BESS agent is a state-of-charge integrator with separate charge/discharge efficiencies $\eta_{\mathrm{ch}},\eta_{\mathrm{dis}}$, energy capacity $E$, and dispatch interval $\Delta T$,
\begin{equation}
  \mathrm{SoC}_{t+1}=
  \begin{cases}
    \mathrm{SoC}_t-\dfrac{P_t\,\Delta T}{\eta_{\mathrm{dis}}\,E},
      & P_t>0\ \text{(discharge)},\\[6pt]
    \mathrm{SoC}_t+\dfrac{|P_t|\,\eta_{\mathrm{ch}}\,\Delta T}{E},
      & P_t<0\ \text{(charge)},
  \end{cases}
  \label{eq:bess_soc}
\end{equation}
held within a physical SoC band. A V2G agent obeys the same recursion augmented with availability connection at step $t$, a discharge-enable SoC threshold, a usable-capacity cap, and a first-order response lag with time constant $T_{\mathrm{v2g}}$, and leaves the fleet at its departure step. DPV operates at maximum power point with curtailment-only (non-positive) regulation, treated as instantaneous at the fast step.

\subsection{Located Disturbance and Discretization}
Eliminating the passive buses with the same partition as the Kron reduction (7) of the main paper maps a bus-level injection to the retained GFM units through the Kron transfer $J_L=J_{IL}J_{LL}^{-1}$, so the input matrix is topology-dependent,
\begin{equation}
  B_{\mathrm{net}}(\mathcal{G}_t)=\operatorname{diag}\!\big(\tfrac{1}{2H}\big)\,J_L(\mathcal{G}_t).
  \label{eq:bnet}
\end{equation}
A uniform, location-agnostic injection (the implicit assumption of a scalar-COI forcing) discards this dependence and is recovered only in the degenerate case $J_L\!\propto\!\mathbf{1}$. Because $\Delta t=1$\,s is large relative to the inverter time constants, the continuous dynamics (5) of the main paper are discretized by zero-order hold,
\begin{equation}
  x_{t+1}=\Phi\,x_t+\Gamma\,u_t,\qquad
  \Phi=e^{A_f\Delta t},\quad
  \Gamma=\!\int_0^{\Delta t}\! e^{A_f\tau}\,d\tau\,B,
  \label{eq:zoh}
\end{equation}
with $(\Phi,\Gamma)$ evaluated by the augmented matrix exponential
\begin{equation}
  \exp\!\left(
  \begin{bmatrix} A_f & I\\ \mathbf 0 & \mathbf 0\end{bmatrix}\Delta t\right)
  =\begin{bmatrix} \Phi & M\\ \mathbf 0 & I\end{bmatrix},
  \qquad \Gamma = M B,
  \label{eq:augexp}
\end{equation}
which is robust to a singular $A_f$ and avoids the forward-Euler overshoot when $\Delta t$ exceeds the system time constants; $(\Phi,M)$ is cached per (topology, droop-gain bin) pair, the gain bin quantizing the mean $|K|$.

At run time the disturbance forcing is resolved by location, in priority order: (i)~an event at a GFM bus enters that unit's swing directly, $-\tfrac{1}{2H_g}\Delta P_{\mathrm{imb}}$; (ii)~an event at a mapped passive bus enters through the corresponding (one-hot) column of $J_L$ via \eqref{eq:bnet}; (iii)~an explicit per-passive-bus vector enters through \eqref{eq:bnet} in full; and (iv)~an unlocatable event (e.g., a line trip that changes the partition itself) falls back to a rating-share injection $-\operatorname{diag}(\tfrac{1}{2H})\,\mathbf s\,\Delta P_{\mathrm{imb}}$ with $s_g=S_g/\sum_g S_g$. The nadir-projection layer (Sec.~\ref{supp:projection}) uses the \emph{same} resolution, so the guard predicts exactly the forcing the plant applies.

\subsection{Aggregate Inertia and Secondary Loop}
The only lumped aggregate the model carries is the rating-weighted backbone inertia,
\begin{equation}
  H_{\mathrm{sys}}=\sum_g H_g\,S_g/S_{\mathrm{base}},
  \label{eq:hsys}
\end{equation}
where each $H_g$ is programmable virtual inertia, so $H_{\mathrm{sys}}$ is set by control rather than by physical machines. A slow integral secondary loop runs concurrently with the primary response and removes the steady-state offset,
\begin{equation}
  \zeta_{t+1}=\operatorname{sat}\!\big(\zeta_t+k_{\mathrm{agc}}\,\Delta f_t\,\Delta t\big),
  \label{eq:agc}
\end{equation}
with $k_{\mathrm{agc}}=0.05$~pu/Hz/s and anti-windup saturation at $\pm0.35$~pu. The integral state feeds back as a rating-share correction to the reference channel, $\Delta P_{\mathrm{ref}}^{\mathrm{eff}}=\Delta P_{\mathrm{ref}}-\zeta_t\,\mathbf s$, applied with a one-step delay so that no algebraic loop forms with the read-out (8) of the main paper. Running the secondary loop \emph{concurrently} with the primary response (the textbook two-loop design) rather than freezing it while FFR is active is essential for sustained disturbances (e.g., a renewable surplus): a frozen integrator can leave the frequency parked off-nominal whenever the primary layer alone cannot re-enter the FFR-deactivation band.

\subsection{Worst-Unit Transient Diagnostic}
The relative-angle block of the state matrix (6) of the main paper lets the per-unit deviations $\Delta f_{g,t}=f_0\,\Delta\omega_{g,t}$ differ during the transient; these are reported only as a spatial diagnostic through the worst-unit excursion
\begin{equation}
  \Delta f^{\mathrm{wc}}_t=\max_{g}\,\big|\Delta f_{g,t}\big|,
  \qquad
  f^{\mathrm{nadir}}_{\mathrm{wc}}=\min_{g,t}\big(f_0+\Delta f_{g,t}\big),
  \label{eq:worstunit}
\end{equation}
evaluated within the model's validity band ($\pm5$~Hz), not as the headline metric.

\subsection{Realized Control Law and Per-GFM Coupling}
\label{supp:coupling}
The two action channels of each agent combine into the realized per-DER injection $\Delta P^{\mathrm{ffr}}_{m,t}=\Delta P_m-K_m\,\Delta f_t$ (main paper, Sec.~III-A). The actuator path applies, in order: a per-channel slew limit $|a_t-a_{t-1}|\le0.2$ per fast step ($\le20\,\%$ of rating per second, which keeps step-like setpoint changes inside the $2$~Hz/s ride-through bound while still allowing a full ramp within $\sim\!5$~s); a clip to the device rating and SoC headroom; and the FFR-eligibility gate $K_m=0$ whenever $\mathrm{SoC}_m\notin[0.25,0.85]$ (nested with $0.05$ margin inside the physical band $[0.20,0.90]$, so a committed DER can always deliver a bidirectional response without striking its energy limits).

For the frequency model, the per-DER commands aggregate to the electrically nearest GFM unit $g(m)$, the same mapping the observation builder uses. The reference channel adds in per unit on the system base,
\begin{equation}
  \Delta P_{\mathrm{ref},g}
   = \frac{1}{S_{\mathrm{base}}}\sum_{m:\,g(m)=g}\Delta P_m ,
  \label{eq:pref_agg}
\end{equation}
while the learned droop gains, which are absolute quantities in MW/Hz, enter the dimensionless per-unit damping of the swing model after unit conversion,
\begin{equation}
  K_{\mathrm{eff},g}
   = K_{\mathrm{bb}}
   + \frac{f_0}{S_{\mathrm{base}}}\sum_{m:\,g(m)=g} K_m ,
  \label{eq:kdroop_conv}
\end{equation}
where $K_{\mathrm{bb}}$ is the backbone droop. The conversion follows from $\Delta P_{\mathrm{MW}}=-K_m\,\Delta f$, $\Delta f=f_0\,\Delta\omega_{\mathrm{pu}}$, and $\Delta P_{\mathrm{pu}}=\Delta P_{\mathrm{MW}}/S_{\mathrm{base}}$; parallel droop sources add, so the learned grid-following droop \emph{adds to} the backbone damping in $D(K_t)=\operatorname{diag}(K_{\mathrm{eff},g})$ of (6), preserving the steady-state droop identity on the per-unit scale of the model.

% ==================================================================
\section{Stability Certificate for the Topology-Aware Model}
\label{supp:lyapunov}
% ==================================================================

This section gives the modal structure of the closed loop and then certifies that the system matrix $\mathbf{A}_f(G,\mathbf{K})$, (6) of the main paper, is stable over the \emph{entire} feasible set, not merely at a nominal operating point.

\subsection{Modal Structure of the Closed Loop}
\label{supp:modal}
Angle-reference invariance makes the rows of the lossless part of $\tilde J_r$ sum to zero; with the resistive losses of the distribution feeder retained, the row sums are small but nonzero. Projecting the swing block onto the rating weights $s_g=S_g/\sum_g S_g$ therefore decouples the COI to leading order: the rating-weighted average is governed by a \emph{single real mode}
\begin{equation}
  \lambda_{\mathrm{COI}}
   \;\approx\;
   -\,\frac{\sum_g K_{\mathrm{eff},g}}{\sum_g 2H_g},
  \label{eq:coi_mode}
\end{equation}
($\tau\!=\!1.5$~s at the backbone gains of Table~I of the main paper), while the remaining $2(n_g-1)$ eigenvalues form inter-machine swing pairs at
$\omega_i\approx\sqrt{\omega_0\,\lambda_i(\tilde J_r)/2\bar H}$. Because the Kron-reduced synchronizing coefficients of the short, stiff feeder are of order $10^3$--$10^4$~pu, these pairs lie between $\approx\!38$ and $\approx\!345$~Hz for the base topology, with decay rates $\operatorname{Re}\lambda\in[-1.0,-0.36]$~s$^{-1}$ ($\tau\le2.8$~s), two orders of magnitude above the $1$~s control bandwidth, and they cancel in the rating-weighted read-out (8). Three consequences follow: (i)~the COI trajectory is quasi-first-order, with none of the governor-style oscillatory tail of synchronous-machine SFR models; (ii)~the learned action shapes the recovery through $D(K_t)$ and $\Delta P_{\mathrm{ref}}$ rather than by phase-shaping an oscillatory mode; and (iii)~the COI-level guard of Sec.~\ref{supp:projection} acts on the only mode that is resolvable and actionable at the control step. The inter-machine pairs are excited by located disturbances but remain inside the model's small-signal validity band and are reported only through the worst-unit diagnostic \eqref{eq:worstunit}.

\subsection{Certificates over the Topology--Gain Set}
\emph{Frozen-time stability.} Because $\mathbf{A}_f$ is affine in the effective gains $K_{\mathrm{eff},i}$ through the damping block $D(K_t)$, the LMI residual $\mathbf{A}_f^{\!\top}\mathbf{P}+\mathbf{P}\,\mathbf{A}_f$ is affine in $K$ for any fixed $\mathbf{P}\succ\mathbf 0$. Verifying the common-$\mathbf{P}$ condition \eqref{eq:clf_lmi} at the $2^{n_K}$ box vertices therefore certifies it over the full box $[0,K_i^{\max}]$ by convex combination, since existence of a common quadratic Lyapunov function at all vertices is equivalent to its existence over the entire gain polytope. For each feasible topology $G\in\mathcal{G}^{\mathrm{feas}}$ and each vertex $\mathbf{K}$ of the box we form $\mathbf{A}_f(G,\mathbf{K})$ and verify Hurwitzness. Over the held topology set this yields a finite family $\{\mathbf{A}_f^{(q)}\}$; every member is Hurwitz (Sec.~V-B of the main paper). The per-vertex Hurwitz test alone is necessary but, since the set of Hurwitz matrices is non-convex, not sufficient for box-wide stability without the shared $\mathbf{P}$.

\emph{Switching stability.} Frozen-time stability is necessary but not sufficient under switching. We first test for a common Lyapunov function (CLF) by the linear matrix inequality
\begin{equation}
\mathbf{P}\succ\mathbf 0,\quad
\mathbf{A}_f^{(q)\top}\mathbf{P}+\mathbf{P}\,\mathbf{A}_f^{(q)}
\preceq-\varepsilon\mathbf{I}\;\;\forall q,
\label{eq:clf_lmi}
\end{equation}
solved as a semidefinite program over the vertex family $\{\mathbf{A}_f^{(q)}\}$. The closed loop switches on two timescales: the gains $K_t$ may change every $1$~s control step, whereas topology $G$ changes only on DSO reconfiguration (minutes). When \eqref{eq:clf_lmi} is feasible, $V(x)=x^{\!\top}\mathbf{P}x$ is a common Lyapunov function for the whole family; by the affine-in-$K$ argument above it remains one for every interior gain, so step-wise $K_t$ changes and topology changes are both admissible \emph{arbitrary} switches and exponential stability holds with no dwell-time condition on either. When \eqref{eq:clf_lmi} is infeasible, as expected when the topology family is diverse, we retain a per-topology $\mathbf{P}_G$ certifying the gain box \emph{within} that fixed $G$ (the affine-in-$K$ vertex argument applies for a fixed topology, so continuous $K_t$ variation inside a mode is still covered), and stability across topology modes is certified by an \emph{average dwell-time} argument: per-mode Lyapunov functions give a switching threshold $\tau_a^{\star}=\ln\mu/(2\eta)$, where $\eta$ is the slowest strictly-damped decay rate and $\mu$ bounds the inter-mode Lyapunov conditioning $\lambda_{\max}(\mathbf{P}_{G'})/\lambda_{\min}(\mathbf{P}_G)$. Concretely, the held family is enumerated exhaustively rather than sampled: $|\mathcal{G}^{\mathrm{feas}}|=24$ feasible topologies $\times\,2^{n_{\mathrm{gfm}}}=4$ gain-box vertices ($n_{\mathrm{gfm}}=2$ GFM units, effective gain on $[K_{\mathrm{bb}},K_i^{\max}]$) give $96$ frozen-time matrices, all Hurwitz ($96/96$) with worst spectral abscissa $\lambda_{\max}=-7.29\times10^{-6}\,\mathrm{s^{-1}}$. The common-$\mathbf{P}$ LMI \eqref{eq:clf_lmi} is infeasible over this family, so the average dwell-time branch governs. With slowest strictly-damped rate $\eta=0.283\,\mathrm{s^{-1}}$ and inter-mode conditioning $\mu=1.69\times10^{10}$, the threshold is $\tau_a^{\star}=\ln\mu/(2\eta)=41.6\,\mathrm{s}$. Since DSO reconfiguration is event-triggered at the planning timescale (minutes), the realized dwell-time exceeds $\tau_a^{\star}$ by orders of magnitude, so the switched system is exponentially stable. The single near-marginal common-mode (COI) frequency left by the relative-angle formulation is excluded from the strictly-damped subspace and is stabilized by the secondary loop~\eqref{eq:agc}. The numerical certificate (vertex Hurwitz check, CLF test, and dwell-time) is produced by the accompanying \texttt{lyapunov\_certificate} routine.

% ==================================================================
\section{Policy-Optimization Loss}
\label{supp:ppo}
% ==================================================================
With the importance ratio $\rho_t$ and the GAE advantage $\hat A_t$ of the main paper, the parameter-shared actor is updated by the clipped PPO surrogate
\begin{equation}
  \mathcal{L}^{\mathrm{clip}}
   = -\,\mathbb{E}\!\left[
       \min\!\big(\rho_t\hat{A}_t,\
       \operatorname{clip}(\rho_t,1-\epsilon,1+\epsilon)\hat{A}_t\big)\right],
  \label{eq:ppo_clip}
\end{equation}
with clip range $\epsilon=0.2$; the total loss adds a value-regression term (coefficient $0.5$) and an entropy bonus annealed across the training curriculum. The action executed in the environment and stored in the rollout is the clipped variate $a_t=\operatorname{clip}(\tilde a_t,-1,1)$ with $\tilde a_t\sim\mathcal{N}(\mu_t,\sigma_t^2)$; both the behaviour log-probability $\log\pi_{\theta_{\mathrm{old}}}(a_t)$ and the target $\log\pi_\theta(a_t)$ entering $\rho_t$ are evaluated at this same $a_t$ under the Gaussian density, so the importance ratio is consistent with the executed action. The density-at-boundary approximation incurred when $|\tilde a_t|>1$ is bounded by the empirical saturation rate ($\le2\,\%$).

% ==================================================================
\section{Closed-Form Nadir Projection}
\label{supp:projection}
% ==================================================================
At every fast step, before the dynamics are integrated, the environment applies a minimal-perturbation projection to the per-GFM reference vector $\Delta P_{\mathrm{ref}}\in\mathbb{R}^{n_g}$. With the cached ZOH pair $(\Phi,M)$ of \eqref{eq:augexp}, the $\omega$-block row selection $(\cdot)_\omega$, the COI weights $w_g=S_g/\sum_g S_g$, and the input gain $G=\operatorname{diag}(\tfrac{1}{2H})$, the next-step COI deviation is \emph{affine} in the reference,
\begin{equation}
  \Delta f^{+}=a+b^{\!\top}\Delta P_{\mathrm{ref}},
  \label{eq:affine_pred}
\end{equation}
\begin{equation}
  a=f_0\,w^{\!\top}\!\big[(\Phi x_t)_\omega+M_{\omega\omega}\,c_u\big],
  \qquad
  b=f_0\,G\,M_{\omega\omega}^{\!\top}\,w,
  \label{eq:affine_coeffs}
\end{equation}
where $M_{\omega\omega}$ is the $\omega$-block of $M$ and $c_u$ collects the reference-independent inputs: the one-step-delayed AGC feedback $-G\,\mathbf s\,\zeta_t$ and the located disturbance forcing resolved exactly as in Sec.~\ref{supp:models} (cases i--iv). If $\Delta f^{+}\in[-\delta,\delta]$ the action passes unchanged; otherwise it is Euclidean-projected onto the violated half-space,
\begin{equation}
  \Delta P_{\mathrm{ref}}\;\leftarrow\;
  \Delta P_{\mathrm{ref}}
  + b\,\frac{\mp\delta-\Delta f^{+}}{\lVert b\rVert^{2}},
  \label{eq:projection}
\end{equation}
the closed-form solution of $\min\lVert\Delta P'-\Delta P_{\mathrm{ref}}\rVert^2$ subject to $a+b^{\!\top}\Delta P'=\mp\delta$. One linear constraint means no quadratic program is solved in the loop; because the predictor reuses the cached $(\Phi,M)$ and the shared forcing resolution, guard and plant evolve under \emph{identical} dynamics, and the layer is applied identically to every learning controller (main paper, Sec.~III-D).

% ==================================================================

% ==================================================================
\section{Training Protocol: Curriculum and Hyperparameter Search}
\label{supp:training}
% ==================================================================
Training follows a six-phase disturbance curriculum (Table~\ref{tab:curriculum}) that escalates the event probability, the maximum imbalance, and the event mix from load-only steps to the full $N{-}1$ mix at $40\,\%$ of $S_{\mathrm{base}}$, while annealing the learning rate and the entropy bonus toward pure exploitation in the final phase. All learning baselines are trained under the identical schedule and budget.

\begin{table}[!t]
\caption{Six-phase disturbance curriculum ($6000$ episodes).}
\label{tab:curriculum}
\centering
\footnotesize
\setlength{\tabcolsep}{3.5pt}
\begin{tabular}{lcccccc}
\toprule
Phase & Episodes & $p_{\mathrm{event}}$ & $|\Delta P|_{\max}$ & Event mix (ls/gt/lt/hr) & LR$\times$ & Entropy\\
\midrule
A & $400$  & $0.50$ & $2.5$~MW ($16\,\%$) & $1/0/0/0$            & $1.0$  & $5\!\times\!10^{-3}$\\
B & $1000$ & $0.70$ & $3.0$~MW ($19\,\%$) & $.5/.5/0/0$          & $1.0$  & $3\!\times\!10^{-3}$\\
C & $1000$ & $0.80$ & $4.0$~MW ($25\,\%$) & $.35/.35/0/.3$       & $1.0$  & $2\!\times\!10^{-3}$\\
D & $1000$ & $0.85$ & $5.5$~MW ($35\,\%$) & $.30/.30/.15/.25$    & $0.8$  & $1\!\times\!10^{-3}$\\
E & $1000$ & $0.90$ & $5.5$~MW ($35\,\%$) & $.3/.3/.2/.2$        & $0.5$  & $5\!\times\!10^{-4}$\\
F & $1600$ & $0.90$ & $6.3$~MW ($40\,\%$) & $.3/.3/.2/.2$        & $0.25$ & $0$\\
\bottomrule
\end{tabular}
\\[2pt]\footnotesize ls = load step, gt = generation trip, lt = line trip, hr = high-renewable surge.
\end{table}

Hyperparameters are selected by asynchronous successive halving (ASHA): $12$ sampled configurations, elimination factor $\eta=3$, stage budgets of $25/75/225$ episodes, and up to three random seeds per surviving configuration at each promotion, scored by the evaluation ITAE on the training feeder. The selected configuration (Table~V of the main paper) is then trained under the full curriculum; checkpoints from intermediate curriculum phases are never reported.

% ==================================================================
\section{Evaluation Metrics and Audit Methodologies}
\label{supp:metrics}
% ==================================================================

\subsection{Frequency-Security Metrics}
With event time $t_e$, episode length $T$, and COI deviation $\Delta f(t)$, the primary settling metric is the time-weighted integral error
\begin{equation}
  \mathrm{ITAE}=\int_{t_e}^{T}(t-t_e)\,|\Delta f(t)|\,dt,
\end{equation}
which weights the deviation by elapsed time and so penalizes the slow, oscillatory post-event tail of a never-settling recovery. The settling time is
$t_{\mathrm{set}}=\inf\{\tau:\,|\Delta f(t)|\le\varepsilon\ \forall t\ge\tau\}-t_e$ with $\varepsilon=50$~mHz, the ENTSO-E SO~GL standard frequency range. FFR success is the indicator that the \emph{total} time the COI frequency spends below the $49.5$~Hz alarm does not exceed $300$~ms (a conservative UFLS-coordination margin, Sec.~\ref{supp:hierarchy}), evaluated on the $0.1$~s sub-step trace. Nadir, zenith, and maximum RoCoF are reported at the control timescale, i.e., on the $1$~s ZOH samples the controller observes, is rewarded on, and is guarded on.

\subsection{Sub-Step Traces for Figures}
The ZOH step lands on the quasi-steady value of each fast step and does not resolve the sub-second transient. For figures only, each fast step is re-simulated at micro-step $\Delta t/n_{\mathrm{sub}}$ using the group property $e^{A_f\Delta t}=(e^{A_f\Delta t/n_{\mathrm{sub}}})^{n_{\mathrm{sub}}}$ with the identical input vector; the trace's last sample coincides with the propagated state to round-off, the re-simulation is read-only, and the tabulated metrics remain those of the control timescale. Figure annotations always quote the tabulated (control-timescale) extrema so figures and tables cannot disagree.

\subsection{Topology Hold-Out Distance}
Reconfigurations are compared by the Jaccard edge distance
$d_J(G_a,G_b)=1-|E_a\cap E_b|/|E_a\cup E_b|$;
the $24$-topology hold-out spans $d_J\le0.18$ from the training feeder, so the protocol probes structural generalization within the reachable radial configuration space, not transfer to arbitrary graphs.

\subsection{Per-Event DSO Cost}
\label{supp:dsocost}
Each event is settled by (15) of the main paper as a deliverability-weighted procurement cost plus an emergency externality,
\begin{equation}
  C_{\mathrm{DSO}}
   = W_{\mathrm{del}}\big(\lambda^{\mathrm{cap}}R_{\mathrm{armed}}T
     + \lambda^{\mathrm{act}}E_{\mathrm{FFR}}\big)
   + \mathrm{VOLL}\cdot E_{\mathrm{shed}},
  \label{eq:dsocost}
\end{equation}
with the armed reserve $R_{\mathrm{armed}}=K_{\mathrm{peak}}\,\delta$ (committed droop at the band edge), the delivered FFR energy $E_{\mathrm{FFR}}$, and the energy-not-served $E_{\mathrm{shed}}$ of a staged ENTSO-E UFLS model applied to the realized COI trajectory (cumulative shed fractions $5/15/30\,\%$ of load below $49.0/48.5/48.0$~Hz; Sec.~\ref{supp:hierarchy}). The capacity term prices the \emph{committed} reserve procured ex-ante and the activation term the \emph{delivered} energy: $K_{\mathrm{peak}}$ is the peak committed droop over the event, $R_{\mathrm{armed}}$~[MW]$=K_{\mathrm{peak}}\delta$ with $\delta=0.5$~Hz, and $E_{\mathrm{FFR}}$~[MWh] is the integrated delivered FFR energy. Units balance with $T$ in \emph{hours}, $T=N_{\mathrm{step}}\Delta t/3600$ ($\Delta t=1$~s over the $50$~s window, $T=50/3600$~h), so $\lambda^{\mathrm{cap}}$~[\texteuro/MW/h]$\cdot R_{\mathrm{armed}}\cdot T$ and $\lambda^{\mathrm{act}}$~[\texteuro/MWh]$\cdot E_{\mathrm{FFR}}$ are both in euros. When no droop is armed ($\sum_i K_i=0$, e.g.\ the no-FFR control) we set $W_{\mathrm{del}}=1$ by convention and the reserve terms vanish ($R_{\mathrm{armed}}=0$), leaving only the VOLL externality, so \eqref{eq:wdel} never divides by zero. We report a per-event procurement cost under a fixed administrative tariff, not a cleared auction price.

Energy and FFR reserve are \emph{distinct co-optimized products} and are priced as such, which resolves an apparent inconsistency between the locational pricing of Sec.~\ref{supp:hierarchy} and this settlement. In a $100\,\%$-renewable islanded microgrid the marginal \emph{energy} cost is essentially zero: the Layer-0 SOCP duals give an energy component $\lambda_p^{\mathrm{energy}}=0$ and a load-weighted DLMP of order $10^{-4}$~\texteuro/MWh, dominated by its loss/congestion part. The FFR \emph{reserve} product is therefore priced by an administrative ancillary tariff that sets the cost \emph{level} ($\lambda^{\mathrm{cap}}=50$~\texteuro/MW/h, $\lambda^{\mathrm{act}}=100$~\texteuro/MWh). These settlement tariffs are distinct from the time-varying prices the policy observes during training (capacity drawn from $[8,20]$~\texteuro/MW/h, energy up to $100$~\texteuro/MWh), which the market reward normalizes via $\lambda^{\mathrm{cap}}_{\mathrm{ref}}=20$ and $\lambda^{\mathrm{lmp}}_{\mathrm{ref}}=100$. The locational congestion component $\lambda_p^{\mathrm{cong}}$ of the DLMP then sets the deliverability \emph{shape} through a \emph{heuristic locational deliverability weighting}
\begin{equation}
  W_{\mathrm{del}}=\frac{\sum_i w_i\,K_i}{\sum_i K_i},\qquad
  w_i=\frac{|\lambda_p^{\mathrm{cong}}(b_i)|}{\overline{|\lambda_p^{\mathrm{cong}}|}},
  \label{eq:wdel}
\end{equation}
a commitment-weighted factor over the DER buses $b_i$, normalized by the DER-bus mean so the fleet-average controller stays exactly on the tariff and only the \emph{spatial} distribution of committed droop $K_i$ moves cost: a fleet leaning on congested buses pays more, one on well-connected buses less. The magnitude $|\lambda_p^{\mathrm{cong}}|$ is used because deliverability scarcity scales with the \emph{degree} of congestion at a bus irrespective of its sign; $W_{\mathrm{del}}$ is an explicit heuristic surcharge on this scarcity, not a settlement rule derived from a market-clearing model. The VOLL term is load-side and is left unweighted; $\mathrm{VOLL}=3000$~\texteuro/MWh is the reserve-context (operational) convention rather than the higher reliability-planning value. A sweep over $\{1500,3000,8000\}$~\texteuro/MWh (Table~\ref{tab:voll_sweep}) leaves the ranking unchanged \emph{within the standards-compliant set} (the four learning controllers carry zero shedding, so their totals are VOLL-independent), while only the non-compliant fixed-droop and no-FFR costs scale with VOLL; the proposed controller is cheapest at every level.

Empirically the per-DER weights span $w_i\in[0.52,1.73]$, but the realized fleet factor stays within $W_{\mathrm{del}}\in[0.97,1.00]$ across all methods and scenarios: the controllers do not concentrate droop on congested buses, because the co-optimized PSO placement (\texttt{official\_placement\_v3}) already sites the DERs at electrically well-connected buses. Locational deliverability is thus a second-order ($\le3\,\%$) correction here, and the cost ranking is driven by committed-reserve \emph{volume} and by \emph{avoided shedding}: the externality term dominates for the non-responsive and fixed-droop fleets.

\begin{table}[!h]
\caption{Mean per-event DSO cost (\texteuro) over the $24$ unseen topologies under the VOLL sweep. The four learning controllers shed no load, so their cost is VOLL-independent; only the non-compliant controls scale with VOLL.}
\label{tab:voll_sweep}
\centering
\footnotesize
\begin{tabular}{lccc}
\toprule
Method & VOLL $1500$ & VOLL $3000$ & VOLL $8000$ \\
\midrule
GraphSAGE-MAPPO (\emph{ours}) & \textbf{1.98} & \textbf{1.98} & \textbf{1.98} \\
GCNN-PPO                      & 2.20 & 2.20 & 2.20 \\
MATD3                         & 2.26 & 2.26 & 2.26 \\
MLP-MAPPO                     & 4.44 & 4.44 & 4.44 \\
\midrule
Fixed droop$^{\dagger}$       & 4.21 & 6.47 & 14.03 \\
No FFR$^{\dagger}$            & 244.6 & 489.0 & 1304.0 \\
\bottomrule
\end{tabular}
\\[3pt]
\raggedright\footnotesize $^{\dagger}$Not standards-compliant (security band breached); shown for reference.
\end{table}
% --- end VOLL sweep table ---

% ==================================================================
\section{Frequency-Control Hierarchy and Grid-Code Alignment}
\label{supp:hierarchy}
% ==================================================================
The frequency-security quantities used throughout the paper are anchored to the standard reserve cascade, summarized in Table~\ref{tab:hierarchy}. After an active-power imbalance the system frequency is contained and then restored by a sequence of progressively slower layers: the synchronous (here \emph{virtual}) inertial response sets the initial RoCoF; Fast Frequency Reserve (FFR) and primary droop (FCR) contain the deviation within seconds; automatic and manual frequency-restoration reserves (aFRR/mFRR) and replacement reserves (RR) return the frequency to nominal over minutes; and under-frequency load shedding (UFLS) is the last-resort emergency layer that acts only once the faster layers have failed. The learned controller of this work operates in the two fastest \emph{controllable} layers (by construction a topology-aware FFR-plus-primary-droop resource), while the inertial floor is fixed by the grid-forming backbone, the restoration layer is modeled as the concurrent secondary loop~\eqref{eq:agc}, and the emergency layer enters only as the UFLS boundary that defines the worst-case cost (Sec.~\ref{supp:metrics}).

\begin{table}[!t]
\caption{Standard frequency-control cascade and its realization in this work. Activation times are full-activation (worst-case) values.}
\label{tab:hierarchy}
\centering
\footnotesize
\setlength{\tabcolsep}{3pt}
\begin{tabular}{@{}l p{2.45cm} p{3.05cm}@{}}
\toprule
Layer & Standard trigger / time & Realization here\\
\midrule
Inertial & instantaneous; $\mathrm{RoCoF}=\Delta P/2H_{\mathrm{sys}}$ (IEEE~1547-2018 Cat~III ride-through; ENTSO-E CE $1$~Hz/s) & virtual inertia $H_g$, \eqref{eq:hsys}; non-actionable by the GFL action\\
\addlinespace
FFR & $0.7$--$1.3$~s to full output at $49.5$/$49.6$/$49.7$~Hz; support $\ge5$/$30$~s (Nordic FFR) & learned $\Delta P^{\mathrm{ffr}}_m=\Delta P_m-K_m\Delta f$ at the $1$~s step, $\le20\,\%$/s slew (Sec.~\ref{supp:coupling})\\
\addlinespace
FCR (primary) & $\le30$~s to full activation; droop, deadband $\le10$~mHz (ENTSO-E SO~GL) & $K_{\mathrm{eff}}=K_{\mathrm{bb}}+{}$learned droop, \eqref{eq:kdroop_conv}; COI mode \eqref{eq:coi_mode}\\
\addlinespace
aFRR (secondary) & $\le5$~min; integral restoration to $\pm50$~mHz (ENTSO-E SO~GL) & concurrent AGC integrator \eqref{eq:agc}, $k_{\mathrm{agc}}=0.05$\\
\addlinespace
mFRR\,/\,RR (tertiary) & $\le12.5$--$15$~min\,/\,$\le30$~min; manual (ENTSO-E EB~GL) & beyond the second-scale episode horizon; not modeled\\
\addlinespace
UFLS (emergency) & staged $49.0\!\to\!48.1$~Hz, ${\sim}8$ steps, ${\sim}100$~ms each (ENTSO-E CE); Stage-1 total operating ${\sim}0.3$~s (NERC PRC-006) & staged VOLL shed $5/15/30\,\%$ below $49.0/48.5/48.0$~Hz (Sec.~\ref{supp:metrics})\\
\bottomrule
\end{tabular}
\end{table}

Three numeric levels recur in the metrics, the reward~\eqref{eq:r_nadir}--\eqref{eq:r_ffr}, and the figure guide lines; they denote \emph{distinct} grid-code events that must not be conflated. (i)~The $\pm50$~mHz band is the ENTSO-E SO~GL \emph{standard frequency range} (Continental Europe) and fixes both the settling tolerance and the neighborhood of the frequency deadband. (ii)~The $49.5$~Hz line is the \emph{alarm}/activation level (the lowest of the three Nordic FFR triggers and the point of full FCR-D activation) and serves as the half-band security margin $\delta=0.5$~Hz in the violation and shaping rewards and in the FFR-success criterion. (iii)~The $49.0$~Hz line is the \emph{first UFLS shed stage} (ENTSO-E CE); only this level carries the value-of-lost-load charge in the per-event DSO cost (Sec.~\ref{supp:metrics}). Consequently, the $300$~ms tolerance in the success indicator is a deliberately conservative UFLS-coordination margin rather than a property of the $49.5$~Hz alarm itself: it equals the NERC PRC-006 Stage-1 total operating time and upper-bounds the ${\sim}100$~ms ENTSO-E CE relay delay that would act at the deeper $49.0$~Hz threshold, so a controller that keeps the $49.5$~Hz alarm clear for all but ${\le}300$~ms provably never reaches the first shed stage. The over-frequency mirror ($50.5$~Hz alarm) applies symmetrically to the high-renewable surge case.

\end{document}